Se quiere comprobar si una oscilación libre puede describirse, durante todo un ensayo, mediante
\[
x''+2\alpha x'+\omega_0^2x=0,
\qquad 0<\alpha<\omega_0,
\]
con coeficientes constantes. El tiempo se mide en segundos y el desplazamiento en centímetros. Se han registrado tres máximos positivos consecutivos; el primero define exactamente \(t=0\). Los datos sintéticos son
\[
\begin{array}{c|cc}
\text{máximo}&t\ (\mathrm s)&x\ (\mathrm{cm})\\\hline
1&0&1.000\pm0.020\\
2&2.00\pm0.02&0.500\pm0.010\\
3&4.10\pm0.02&0.250\pm0.020
\end{array}
\]
Cada \(\pm\) indica una cota de error garantizada en la escala indicada; las cotas se pueden combinar independientemente.

Use sólo los dos primeros máximos para obtener intervalos compatibles de \(\alpha\), de la frecuencia angular amortiguada \(\omega_d\) y de \(\omega_0\), distinguiendo las dos frecuencias. Reserve el tercero para contrastar por separado su instante y su altura con las predicciones del modelo. Decida si existe un modelo de coeficientes constantes compatible con todas las cotas; una discrepancia no debe atribuirse a una causa física concreta sin información adicional. Puede utilizar calculadora.
